\chapter*{Abstract (Arabic)}
\addcontentsline{toc}{chapter}{Abstract (Arabic)}

\begin{center}
\textit{Note for the author: this page is a placeholder for the Arabic abstract required by KFUPM. Replace the content below with the Arabic translation of the English abstract. To typeset Arabic correctly, compile this report with \texttt{xelatex} or \texttt{lualatex} after adding \texttt{\textbackslash usepackage\{polyglossia\}} and \texttt{\textbackslash setotherlanguage\{arabic\}} to the preamble.}
\end{center}

\vspace{1cm}

\noindent\textbf{Title (Arabic placeholder):} \\
\textit{[Arabic translation of: \sh{}: A Unified Research Workspace from Discovery to Draft]}

\vspace{1cm}

\noindent\textbf{Abstract Content (to be translated):} \\
\textit{The Arabic abstract should mirror the English abstract content covering: (1) the problem of fragmented academic tooling and citation hallucination in general-purpose AI assistants; (2) the design of \sh{} as a unified research workspace integrating discovery, reference management, collaborative LaTeX editing, and library-grounded AI; (3) the three-tier architecture with backend, frontend, and Y.js-based collaboration server; (4) the three evaluation studies on citation faithfulness, system performance, and comparative workflow; and (5) the contributions and the deployed system at scholarhub.space.}

\vspace{1cm}

\noindent\textbf{Keywords (Arabic placeholder):} \\
\textit{[Translate: collaborative editing, CRDT, retrieval-augmented generation, academic writing tools, paper discovery]}

\clearpage
